\documentclass[12pt]{report}

\usepackage[margin=1in]{geometry}
\usepackage{amsmath}
\usepackage{amssymb}
\usepackage{booktabs}
\usepackage{array}
\usepackage{enumitem}
\usepackage{xcolor}
\usepackage{hyperref}
\usepackage{fancyhdr}

\hypersetup{
  colorlinks=true,
  linkcolor=blue,
  urlcolor=blue
}

\setlength{\parindent}{0pt}
\setlength{\parskip}{0.55em}
\setlength{\emergencystretch}{2em}
\pagestyle{fancy}
\fancyhf{}
\fancyhead[L]{\footnotesize ENGS 101: Principles of Reinforcement Learning}
\fancyhead[R]{\footnotesize Final Project}
\fancyfoot[C]{\thepage}

\newcommand{\intermediatedate}{November 1, 2026, 23:59 ET}
\newcommand{\finalduedate}{November 17, 2026, 23:59 ET}
\newcommand{\repositoryurl}{https://github.com/cccat6/engs101-cyber-rl}
\newcommand{\arenaurl}{https://arena.codefly.run}

\begin{document}

\noindent
\begin{minipage}{\textwidth}
    \raggedright
    \textbf{ENGS 101 Final Project}\\
    \vspace{0.6em}
    \begin{center}
      {\Large \textbf{Red--Blue Cybersecurity RL Challenge}}
    \end{center}
    Instructor: Prof.\ Peter Chin\\
    \textbf{Intermediate:} \intermediatedate\\
    \textbf{Final:} \finalduedate
\end{minipage}

\vspace{0.5em}

\begin{enumerate}[leftmargin=0.25in,label=$\diamond$,itemsep=0.15em]
  \item Submit one ZIP archive through the course website at
        \href{\arenaurl}{\texttt{arena.codefly.run}}. Name the archive with
        your seven-character student ID.
  \item Submit one agent implementation. The same code must support both the
        Red and Blue roles.
  \item You may use the Python standard library, PyTorch, and NumPy. You may
        not use an existing reinforcement-learning implementation or framework.
  \item The project has two separately recorded stages: an Intermediate
        submission and a Final submission. Both stages use the same upload
        format and interface.
  \item The course server performs the official training and evaluation.
        Local results are for development only.
\end{enumerate}

\section*{1. Assessment Schedule and Grading}

Dartmouth's Fall 2026 classes begin on September 14 and end on November 17.
The Intermediate deadline is the end of Week 7, November 1, 2026 at 23:59
Eastern Time. The Final deadline is the last day of classes, November 17, 2026
at 23:59 Eastern Time. These dates are based on the Registrar's Fall Term 2026
calendar: \href{https://registrar.dartmouth.edu/calendars/academic-institutional-calendars/2026-2027-key-dates/fall-term-2026}{Dartmouth Fall Term 2026 calendar}.

The server preserves each upload as a separate timestamped source snapshot.
Please select the correct stage on the upload page, use your own student ID,
and do not share the submission website.

Each stage is scored out of 100 points. An upload receipt confirms only the
archive format. Both role-training runs, model reload and evaluation must
succeed before stage credit is recorded. Fix a failed run and resubmit.
The role score is compared with the fixed course teacher baseline; a role
``beats'' a baseline when its fixed-evaluation point estimate is strictly above
that baseline.

\begin{center}
\begin{tabular}{>{\raggedright\arraybackslash}p{0.74\textwidth}p{0.16\textwidth}}
\toprule
\textbf{Stage result} & \textbf{Stage score} \\
\midrule
Intermediate: both Red and Blue beat the teacher baseline & 100 \\
Intermediate: exactly one role beats the teacher baseline & 80 \\
Intermediate: submission passes validation but neither role beats the baseline & 60 \\
\midrule
Final: both roles beat the Final teacher challenge & 100 \\
Final: exactly one role beats the Final teacher challenge & 80 \\
Final: neither role beats the Final challenge, but both roles beat the Intermediate teacher baseline & 60 \\
Final: neither role beats the Final challenge, and exactly one role beats the Intermediate baseline & 40 \\
Final: neither role beats either teacher baseline & 20 \\
\bottomrule
\end{tabular}
\end{center}

After the Intermediate deadline, the course team freezes the valid Intermediate
policies and trains a stronger Final teacher challenge from that cohort. The
Final teacher challenge, hidden scenarios, random seeds, training budget and
evaluation budget are fixed before Final submissions are evaluated. The
student-vs-student matrix is a diagnostic leaderboard and does not determine
the stage score.

Final uploads remain closed until the Final teacher challenge is frozen and
released by the course team. Intermediate and Final points are retained
separately even when a newer upload replaces the active leaderboard entry.

\section*{2. Objective}

Design and implement a reinforcement-learning agent for a simulated
cybersecurity environment. Red attempts to discover, compromise, and impact
network assets. Blue attempts to detect activity, deploy decoys, remove access,
and restore hosts.

Your code will be trained separately as Red and Blue. The resulting frozen
policies will be scored relative to the frozen course baseline. They will also
play against the policies submitted by the rest of the class in a separate
cross-play matrix.

\section*{3. Starter repository}

Clone the public repository:

\begin{verbatim}
git clone https://github.com/cccat6/engs101-cyber-rl
cd engs101-cyber-rl
python -m venv .venv
\end{verbatim}

Activate \texttt{.venv} using the command for your operating system, then run:

\begin{verbatim}
python -m pip install -r requirements.txt
\end{verbatim}

The repository contains a small public environment and a complete runnable
example. The student-facing files are:

\begin{center}
\begin{tabular}{p{0.28\textwidth}p{0.64\textwidth}}
\toprule
\textbf{File} & \textbf{Purpose} \\
\midrule
\texttt{student\_agent.py} &
Implement your agent here. This is the main file you will edit. \\
\texttt{run.py} &
Train and test the same agent locally as both Red and Blue. Keep this file in
the submission. \\
\texttt{scenarios/base.yaml} &
The public local environment used by \texttt{run.py}. \\
\bottomrule
\end{tabular}
\end{center}

Python 3.11 or later is required. Install the dependencies from
\texttt{requirements.txt}. The file contains PyTorch, NumPy, PyYAML, and the
local environment package.

The public environment is smaller than the hidden server environment. Do not
hard-code observation dimensions, action counts, host names, or a particular
topology.

\section*{4. Required interface}

Keep the following function:

\begin{verbatim}
def create_agent(config):
    return YourAgent(config)
\end{verbatim}

The returned object must expose:

\begin{verbatim}
agent.model
agent.reset(seed)
agent.act(observation, training)
agent.learn(transition)
\end{verbatim}

\texttt{agent.model} must be a \texttt{torch.nn.Module}. The model may contain
at most 1,000,000 registered parameters per role.
The local runner and the course server perform the same check immediately after
\texttt{create\_agent(config)} returns. A submission that exceeds the limit is
rejected before environment training begins.

The other three members have the following meanings:

\begin{center}
\begin{tabular}{p{0.30\textwidth}p{0.62\textwidth}}
\toprule
\textbf{Member} & \textbf{Purpose} \\
\midrule
\texttt{reset(seed)} &
Reset episode-local state. \\
\texttt{act(observation, training)} &
Return one valid integer action. \\
\texttt{learn(transition)} &
Store the transition and perform any training update. \\
\bottomrule
\end{tabular}
\end{center}

The runner supplies this configuration:

\begin{verbatim}
config.role
config.observation_dim
config.num_actions
config.total_steps
config.seed
config.device
\end{verbatim}

Always construct the model from the supplied dimensions and move it to
\texttt{config.device}. The server normally supplies a CUDA device; local
development may use CPU, CUDA, Intel XPU, or Apple MPS.

\section*{5. Observations and transitions}

The policy observation has two fields:

\begin{verbatim}
observation.features
observation.action_mask
\end{verbatim}

\texttt{features} is a one-dimensional \texttt{float32} NumPy array.
\texttt{action\_mask} is a Boolean array with one entry per action.
\texttt{act} must return an action whose mask entry is true.

During training, \texttt{learn} receives:

\begin{verbatim}
transition.observation
transition.action
transition.reward
transition.next_observation
transition.done
\end{verbatim}

The runner calls \texttt{learn} exactly once per environment step. Evaluation
does not call \texttt{learn}.

\section*{6. Implementation rule}

You must implement the RL algorithm yourself using PyTorch tensor operations,
neural-network modules, and optimizers. Allowed components include
\texttt{torch.nn.Linear}, activation functions, automatic differentiation, and
\texttt{torch.optim}.

Do not import an implementation of DQN, PPO, actor--critic, replay buffers,
training loops, or policies from an RL package. In particular, do not use
Stable-Baselines3, RLlib, CleanRL, Tianshou, TorchRL algorithms, or equivalent
ready-made implementations.

\section*{7. Local run}

From the repository root, run:

\begin{verbatim}
python run.py --steps 1000 --device auto
\end{verbatim}

This single command:

\begin{enumerate}[leftmargin=0.3in]
  \item trains the submitted code as Red;
  \item saves and reloads the Red model;
  \item tests the frozen Red model;
  \item repeats the same process as Blue.
\end{enumerate}

After testing your changes, package only the required source files:
\begin{verbatim}
python package_submission.py YOUR_ID
# If your agent imports a helper module:
python package_submission.py YOUR_ID --include helper.py
\end{verbatim}
Replace \texttt{YOUR\_ID} with your own seven-character Dartmouth ID.

The provided \texttt{SimpleAgent} is a minimal one-layer policy-gradient
example. It runs end to end and updates its PyTorch model, but it is deliberately
small and is not intended to be competitive. You may replace its implementation
while preserving the required interface.

\section*{8. Server training and evaluation}

The course server supplies the official environment-step budget, random seeds,
training scenarios, device, and evaluation episodes. These values are the same
for all submissions.

For each role, the runner uses one frozen, course-owned teacher policy as the
opponent. The public repository contains teacher checkpoints for local work;
the server uses separately trained private checkpoints with the same
interface. The teacher opponent is reset between episodes but is never updated
during student training. Thus every submission receives the same number of
environment steps and the same fixed opponent.

After training, the server saves the weights in \texttt{agent.model}. It then
creates a fresh agent, reloads those weights, and performs frozen inference.

The public and hidden environments use the same interface but may have
different dimensions and network topologies.

\section*{9. Reward and score}

The environment supplies a per-step reward for learning. The leaderboard score
is computed separately from the mission loss over the whole episode.

For an episode of length $T$, let

\[
D = \frac{1}{T}\sum_{t=1}^{T}\text{mission\_loss}_t.
\]

The corresponding role scores are

\[
\text{RedScore}=100D,
\qquad
\text{BlueScore}=100(1-D).
\]

For the stable score, each frozen student role is evaluated against the
opposite frozen teacher role. Let $R_s$ be the raw score of a student Red
policy against teacher Blue, and let $R_b$ be teacher Red's raw score against
teacher Blue. Define the relative Red score as

\[
\text{RelativeRed}=\operatorname{clip}(50+R_s-R_b,0,100).
\]

The relative Blue score is defined analogously from student Blue against
teacher Red. The teacher baseline therefore receives 50 for each role. The
stable score is the average of the relative Red and Blue scores. The server
also reports a 95\% confidence interval. This centering changes only the
aggregate, not the environment reward or any episode score. Training rewards
are not leaderboard points.

\section*{10. Leaderboard}

Each accepted submission produces one frozen Red policy and one frozen Blue
policy. The main standing ranks the stable score relative to the teacher
baseline. The server separately evaluates every submitted Red policy against
every submitted Blue policy and the course baseline.

When a new submission is accepted, the server recomputes baseline-relative
scores, rebuilds the comparison matrix using the saved frozen policies, and
re-ranks all participants. Old policies are re-evaluated, not retrained.
The upload form, submission status, stable standings, confidence intervals,
and Red--Blue comparison matrix appear on the same course webpage. Matrix
scores are relative to the current class population and do not determine the
stable score. Hidden scenarios and trajectories are not released.

\section*{11. Submission}

Upload one ZIP archive through the course website. Its filename must be your
seven-character student ID followed by \texttt{.zip}. The ID may contain ASCII
letters and digits, in either case; the server converts it to lowercase. For
example:

\begin{verbatim}
F001ABC.zip
\end{verbatim}

Submit only the runnable entry points and any Python helper modules used by your
agent. The archive may place these source files at its root:

\begin{verbatim}
F001ABC.zip
|-- run.py
|-- student_agent.py
`-- helper.py                 (if used)
\end{verbatim}

Alternatively, the same contents may be inside one folder:

\begin{verbatim}
F001ABC.zip
`-- F001ABC/
    |-- run.py
    |-- student_agent.py
    `-- helper.py             (if used)
\end{verbatim}

If the archive contains a single outer folder, that folder must have the same
student ID as the ZIP filename, ignoring case. After normalization,
\texttt{run.py} and \texttt{student\_agent.py} must be at the submission root.
Every submitted file must be Python source with a \texttt{.py} extension. The
ZIP file and its extracted Python source must each be no larger than 100 KiB.
An archive that does not follow these rules may be rejected or removed from the
submission system. Submit a corrected archive using your own seven-character
student ID in the filename.
Do not include the starter repository, scenarios, dependency files,
\texttt{.venv}, local checkpoints, pretrained weights, or other generated
files. A malformed or oversized archive is rejected by the website with an
error message.

Do not submit pretrained weights. The server trains a fresh Red model and a
fresh Blue model from your code. \texttt{run.py} remains in the archive so the
submission is directly runnable locally. Official training calls
\texttt{create\_agent} in \texttt{student\_agent.py}.

Submissions must:

\begin{enumerate}[leftmargin=0.3in]
  \item import successfully in the fixed course environment;
  \item follow the required interface;
  \item return only valid actions;
  \item use only the permitted libraries;
  \item keep each role's model at or below 1,000,000 registered parameters;
  \item complete within the server resource limits.
\end{enumerate}

The official server training budget and evaluation episode count are controlled
by the course evaluator and are shown on the upload page. The current budget
is 10,000 environment steps per role and five evaluation episodes per hidden
scenario. Deadline checks use the receipt time, not when queued grading finishes.
Both deadlines are in \texttt{America/New\_York}; the UTC offset is $-05{:}00$
on November 1 and November 17. The two stage rules
and deadlines above are the published project grading policy.

\end{document}
